\chapter{Solutions: Midterm, 6 May 2014}
\label{s:theory-2014-05-midterm}
\begin{sourceNote}
Only printed Questions 1--3 survive, totaling 60\%. They are fully solved here. The final section contains supplementary exercises newly written from topics recognizable in handwritten notes; they are not reconstructions of the missing printed questions. The two-dimensional graphite model in the source is a single honeycomb layer, commonly called graphene.
\end{sourceNote}

\section*{1. True or false, with assumptions}
\subsection*{(a) Identical atoms must be related by a Bravais translation: false}
A Bravais translation maps the entire crystal onto itself. Two identical atoms at different basis sites need not be connected by such a translation. Question 3 supplies a counterexample: both basis atoms are carbon, but their separation is
\[
\vct d_2-\vct d_1=a(0,1/\sqrt3).
\]
If this were $n_1\vct t_1+n_2\vct t_2$ with integers $n_i$, the $x$ component would give $n_1=n_2=n$. The $y$ component would require $\sqrt3an=a/\sqrt3$, or $n=1/3$, not an integer.

\subsection*{(b) An exterior point is closer to another reciprocal point: true}
The closed first Brillouin zone is
\[
\mathcal B_1=\{\vct k:|\vct k|\leq|\vct k-\vct G|
\text{ for every }\vct G\in\mathcal L^*\}.
\]
If $\vct k\notin\mathcal B_1$, at least one nonzero $\vct G$ violates this inequality: $|\vct k-\vct G|<|\vct k|$. The assertion follows directly from the Wigner--Seitz definition. A boundary point can be equidistant; being strictly outside is different from being on the boundary.

\section*{2. FCC plane geometry}
The primitive volume and reciprocal basis are
\[
\Omega=\frac{a^3}{4},\qquad
\vct g_1=\frac{2\pi}{a}(-1,1,1),\quad
\vct g_2=\frac{2\pi}{a}(1,-1,1),\quad
\vct g_3=\frac{2\pi}{a}(1,1,-1).
\]
Direct substitution verifies $\vct t_i\cdot\vct g_j=2\pi\delta_{ij}$. One lattice point per primitive cell gives volume density $4/a^3$.

\subsection*{(a) Normal $\tfrac12\vct g_1+\vct g_2-\vct g_3$}
The given vector specifies a direction but is not itself reciprocal. Multiplying by two gives the shortest parallel reciprocal vector:
\[
\vct G_0=\vct g_1+2\vct g_2-2\vct g_3
=\frac{2\pi}{a}(-1,-3,5).
\]
For $\vct R=\sum_i n_i\vct t_i$,
\[
\vct G_0\cdot\vct R=2\pi(n_1+2n_2-2n_3).
\]
The integer on the right can take every value $m$, for instance by setting $n_1=m$, $n_2=n_3=0$. Consecutive occupied planes are therefore $\vct G_0\cdot\vct r=2\pi m$, with
\[
\boxed{d=\frac{2\pi}{|\vct G_0|}=\frac a{\sqrt{35}}.}
\]
Applying $2\pi/|\vct G|$ to the printed half-length normal would skip every other occupied plane.

\subsection*{(b) Surface density}
The integer solutions of $n_1+2n_2-2n_3=0$ are generated by
\[
\vct u=-2\vct t_1+\vct t_2,\qquad
\vct v=2\vct t_1+\vct t_3.
\]
Any in-plane point is $n_2\vct u+n_3\vct v$, so these form a primitive planar cell, not a supercell. Its area is
\[
A_2=|\vct u\times\vct v|=\frac{a^2\sqrt{35}}4,\qquad
\boxed{\rho_2=\frac1{A_2}=\frac4{a^2\sqrt{35}}.}
\]
Independently, $A_2d=\Omega$, giving $\rho_2=(4/a^3)d$. The units are inverse area.

\subsection*{(c) Normal $\vct t_1+2\vct t_2-2\vct t_3$}
Its Cartesian direction is $a(0,-1,3)/2$. Set a parallel reciprocal vector equal to
\[
\vct G=\frac{2\pi}{a}(0,-\lambda,3\lambda)
=h_1\vct g_1+h_2\vct g_2+h_3\vct g_3.
\]
Solving gives $(h_1,h_2,h_3)=(\lambda,3\lambda/2,-\lambda/2)$. The smallest positive $\lambda$ making all coefficients integral is $2$:
\[
\vct G_0=\frac{2\pi}{a}(0,-2,6)=2\vct g_1+3\vct g_2-\vct g_3,
\qquad \boxed{d=\frac a{\sqrt{40}}.}
\]
Since $\gcd(2,3,-1)=1$, every integer plane index occurs and the planes are successive occupied ones.

\section*{3. The two honeycomb bands}
\subsection*{(a) Check differentiability before classifying}
Define
\[
x=\frac{ak_x}{2},\quad y=\frac{\sqrt3ak_y}{2},\quad
S=1+4\cos x\cos y+4\cos^2x,\quad E_\pm=E_p\pm\gamma\sqrt S.
\]
For $S>0$,
\[
\nabla E_\pm=\pm\frac{\gamma}{2\sqrt S}\nabla S.
\]
At $S=0$, this formula cannot simply be evaluated; expand the original dispersion.

\paragraph{$\Gamma=(0,0)$: maximum and minimum.}
Expanding the cosines gives
\[
S=9-\frac32a^2q^2+O(q^4),\qquad
E_\pm=E_p\pm3\gamma\mp\frac{\gamma a^2}{4}q^2+O(q^4).
\]
The gradient vanishes and the Hessian is $\mp\gamma a^2I/2$. Thus the upper band has a maximum $E_p+3\gamma$, and the lower band a minimum $E_p-3\gamma$.

\paragraph{$Q=(\pi/a,\pi/(\sqrt3a))$: saddle.}
Write $\vct k=\vct Q+\vct q$, $u=aq_x/2$, $v=\sqrt3aq_y/2$. Since $x=\pi/2+u$ and $y=\pi/2+v$,
\[
S=1+4u^2+4uv+O(q^4),\qquad
\sqrt S=1+2u^2+2uv+O(q^4).
\]
Consequently
\[
E_\pm=E_p\pm\gamma
\pm\frac{\gamma a^2}{2}(q_x^2+\sqrt3q_xq_y)+O(q^4).
\]
The upper-band Hessian is
\[
\gamma a^2\begin{pmatrix}1&\sqrt3/2\\ \sqrt3/2&0\end{pmatrix},
\]
with eigenvalues $3\gamma a^2/2$ and $-\gamma a^2/2$. Both bands are saddles. In rotated principal-axis coordinates,
\[
E_+-(E_p+\gamma)=\alpha q_1^2-\eta q_2^2+O(q^4),\qquad
\alpha=\frac{3\gamma a^2}{4},\quad\eta=\frac{\gamma a^2}{4}.
\]
The lower band reverses both signs and remains a saddle.

\paragraph{$P=(4\pi/(3a),0)$: Dirac touching.}
Here $S=0$. Expanding gives
\[
S=\frac34a^2q^2+O(q^3),\qquad
E_\pm=E_p\pm\frac{\sqrt3\gamma a}{2}|\vct q|+O(q^2).
\]
Every line through the apex has an absolute-value cusp. The band is therefore \textbf{not differentiable} at the touching. If a critical point means a differentiable point with $\nabla E=0$, $P$ is not one. Under a broader definition allowing nonsmooth extrema, the upper and lower bands have conical minimum and maximum, respectively. One must not infer $\nabla E=0$ from $\nabla S=0$ at $S=0$.

\begin{center}
\begin{tikzpicture}[scale=4,>=Stealth]
\draw[thick,fill=blue!4] (0.666667,0)--(0.333333,0.57735)--(-0.333333,0.57735)--(-0.666667,0)--(-0.333333,-0.57735)--(0.333333,-0.57735)--cycle;
\draw[->] (-0.85,0)--(0.94,0) node[right] {$k_x$};
\draw[->] (0,-0.7)--(0,0.74) node[above] {$k_y$};
\fill (0,0) circle (0.012) node[below left] {$\Gamma$};
\fill[blue] (0.666667,0) circle (0.012) node[above right] {$P$};
\fill[red!70!black] (0.5,0.288675) circle (0.012) node[above right] {$Q$};
\draw[dashed] (0,0)--(0.5,0.288675);
\end{tikzpicture}

\small First Brillouin zone, in units of $2\pi/a$. The source's $P,Q$ are commonly called $K,M$.
\end{center}

\subsection*{(b) There is no gap}
The band separation is $E_+-E_-=2\gamma|F_{\vct k}|\geq0$ and vanishes at $P$ and symmetry-equivalent corners:
\[
\boxed{E_g=0.}
\]
At half filling with $E_F=E_p$, the Fermi surface reduces to Dirac points in this model.

\subsection*{(c) Choose the lower-band minimum at $\Gamma$}
Use DOS per area, per spin, for one local patch:
\[
D(E)=\int_{\rm patch}\frac{\dd^2q}{(2\pi)^2}\delta(E-E(\vct q)).
\]
Let $E_0=E_p-3\gamma$, $A=\gamma a^2/4$, and approximate $E=E_0+Aq^2$ in a circular patch $q<\Lambda$. Setting $z=Aq^2$ gives $q\dd q=\dd z/(2A)$:
\begin{align*}
D_\Gamma(E)&=\frac1{2\pi}\int_0^\Lambda q\dd q\,\delta(E-E_0-Aq^2)\\
&=\boxed{\frac1{\pi\gamma a^2}
\Theta(E-E_0)\Theta(E_0+A\Lambda^2-E).}
\end{align*}
The local approximation gives a constant step at the band edge; the upper cutoff depends on the patch. Multiply by two for both spins and by sample area $\mathcal A$ for a total DOS.

\paragraph{Additional check: the saddle at $Q$.}
For the upper band set $\varepsilon=E-(E_p+\gamma)$. With $\varepsilon=\alpha q_1^2-\eta q_2^2$, use $u=\sqrt\alpha q_1$, $v=\sqrt\eta q_2$, then $s=u+v$, $t=u-v$. Thus $u^2-v^2=st$ and $\dd u\dd v=\dd s\dd t/2$. For a local cutoff $|s|,|t|<L$,
\begin{align*}
D_Q(E)&\simeq\frac1{8\pi^2\sqrt{\alpha\eta}}
\int_{-L}^L\dd s\int_{-L}^L\dd t\,\delta(\varepsilon-st)\\
&=\frac1{4\pi^2\sqrt{\alpha\eta}}\ln\frac{L^2}{|\varepsilon|}\\
&=\frac1{\pi^2\sqrt3\gamma a^2}
\ln\frac{E_c}{|E-(E_p+\gamma)|}.
\end{align*}
This logarithmic van Hove divergence is per saddle and per spin. Its cutoff $E_c$ cannot be fixed by the quadratic expansion alone. In contrast, one Dirac valley gives
\[
D_P(E)=\frac{|E-E_p|}{2\pi\beta^2},\qquad\beta=\sqrt3\gamma a/2,
\]
which vanishes at the touching. These are three distinct DOS signatures.

\clearpage
\section*{Supplementary exercises from handwritten topics}
\begin{sourceNote}
The topics below are recognizable in attached handwritten answers, but their printed question statements, weights, and exact requested scope are missing. Exercises A--D are newly composed learning material, not verbatim reconstructions of original Questions 4--7.
\end{sourceNote}

\subsection*{A. Derive Bloch form from translation symmetry}
\textbf{Exercise.} For a Hamiltonian with Bravais translation symmetry, derive the quasiperiodicity of a simultaneous energy and translation eigenstate.

Define $(T_{\vct R}\psi)(\vct r)=\psi(\vct r+\vct R)$. Then $[H,T_{\vct R}]=0$ and $T_{\vct R}T_{\vct R'}=T_{\vct R+\vct R'}$. Choose a simultaneous eigenstate within each energy eigenspace:
\[
T_{\vct R}\psi=\lambda(\vct R)\psi.
\]
Translations are unitary, so $|\lambda|=1$. The group law imposes $\lambda(\vct R+\vct R')=\lambda(\vct R)\lambda(\vct R')$. Choose the primitive phases as $\lambda(\vct t_i)=\ee^{\ii\vct k\cdot\vct t_i}$; then
\[
\psi(\vct r+\vct R)=\ee^{\ii\vct k\cdot\vct R}\psi(\vct r).
\]
Defining $u_{\vct k}=\ee^{-\ii\vct k\cdot\vct r}\psi$ gives $u(\vct r+\vct R)=u(\vct r)$. Vectors $\vct k$ and $\vct k+\vct G$ have the same translation eigenvalues, so $\vct k$ may be chosen in the first Brillouin zone.

\subsection*{B. Brillouin's theorem and single excitations}
\textbf{Exercise.} Show that a self-consistent HF determinant has zero Hamiltonian matrix element with every singly excited determinant.

Let $\Phi_0$ be the stationary HF determinant and $\Phi_i^a$ replace occupied spin orbital $i$ by virtual orbital $a$. The Slater--Condon rules give
\[
\bra{\Phi_i^a}H\ket{\Phi_0}
=h_{ai}+\sum_{j\in\mathrm{occ}}
\bigl(\langle aj|v|ij\rangle-\langle aj|v|ji\rangle\bigr)=F_{ai}.
\]
For canonical orbitals, $F\chi_i=\epsilon_i\chi_i$ and $\braket{\chi_a}{\chi_i}=0$, hence
\[
\boxed{\bra{\Phi_i^a}H\ket{\Phi_0}=0.}
\]
Equivalently, varying the real and imaginary components of the orbital rotation $\chi_i\mapsto\chi_i+\theta\chi_a$ must give zero first-order energy change. The theorem removes only the coupling to single excitations; double-excitation couplings generally survive, so HF need not be exact.

\subsection*{C. A determinant is an exact noninteracting eigenstate}
\textbf{Exercise.} If $H_0=\sum_{j=1}^Nh(j)$ and $h\chi_i=\epsilon_i\chi_i$, show that the determinant has energy $\sum_i\epsilon_i$.

Expand the determinant:
\[
\Phi=\frac1{\sqrt{N!}}\sum_P(-1)^P\prod_{j=1}^N\chi_{P(j)}(x_j).
\]
For every permutation, $H_0$ acting on the product gives the same eigenvalue $\sum_j\epsilon_{P(j)}=\sum_i\epsilon_i$. Therefore
\[
H_0\Phi=\left(\sum_{i=1}^N\epsilon_i\right)\Phi.
\]
The ground state fills the lowest available $N$ spin orbitals. Adding a true two-body interaction invalidates this argument, which is why an exact noninteracting determinant must be distinguished from an interacting HF approximation.

\subsection*{D. Hartree, Hartree--Fock, and Kohn--Sham}
\textbf{Exercise.} Compare the single-particle equations, variational objects, and treatment of exchange and correlation.

Let $h=-\hbar^2\nabla^2/(2m_e)+v_{\rm ext}$ and $\kappa=e^2/(4\pi\epsilon_0)$.
\begin{enumerate}[label=(\roman*)]
\item \textbf{Hartree.} Vary a product $\prod_i\chi_i$:
\[
\left[h+\sum_{j\ne i}J_j\right]\chi_i=\epsilon_i\chi_i,\qquad
J_j(\vct r)=\int\frac{\kappa|\chi_j(x')|^2}{|\vct r-\vct r'|}\dd x'.
\]
It includes the other electrons' average Coulomb field but no antisymmetry-driven exchange.
\item \textbf{Hartree--Fock.} Vary a single determinant:
\[
\left[h+\sum_j(J_j-K_j)\right]\chi_i=\epsilon_i\chi_i,
\]
\[
(K_j\chi_i)(x)=\chi_j(x)\int
\frac{\kappa\chi_j^*(x')\chi_i(x')}{|\vct r-\vct r'|}\dd x'.
\]
Exchange is nonlocal and $J_i-K_i$ cancels self-interaction. Correlation beyond a single determinant is generally missing.
\item \textbf{Kohn--Sham.} Introduce auxiliary noninteracting orbitals reproducing the interacting density:
\[
F[n]=T_s[n]+E_H[n]+E_{xc}[n],\quad
E_H=\frac\kappa2\iint\frac{n(\vct r)n(\vct r')}{|\vct r-\vct r'|}\dd^3r\dd^3r'.
\]
Orbital variation gives
\[
[h+v_H(\vct r)+v_{xc}(\vct r)]\varphi_i=\epsilon_i\varphi_i,\qquad
n=\sum_i f_i|\varphi_i|^2,
\]
with spin summation understood, $v_H=\delta E_H/\delta n$, and $v_{xc}=\delta E_{xc}/\delta n$. The functional $E_{xc}$ contains both nonclassical interactions and the difference between true and auxiliary kinetic energies. It is not merely HF exchange renamed. An exact functional gives exact ground-state density and energy in principle; approximate functionals add their own errors.
\end{enumerate}

\section*{Primary references used for checks}
\href{https://journals.aps.org/pr/abstract/10.1103/PhysRev.136.B864}{P. Hohenberg and W. Kohn, Inhomogeneous Electron Gas, Phys. Rev. 136, B864 (1964)};
\href{https://journals.aps.org/pr/abstract/10.1103/PhysRev.140.A1133}{W. Kohn and L. J. Sham, Self-Consistent Equations Including Exchange and Correlation Effects, Phys. Rev. 140, A1133 (1965)}.
Exercises A--D are explicitly new exercises based on the note topics. The original printed continuation remains missing.
